% ECONOMICS_PROBLEM: {"id": "F13-409", "title": "Long-run reshoring effects", "jel_code": "F13", "source_id": "OP-0343"}
% !TeX program = xelatex
\documentclass[11pt,a4paper]{article}
\usepackage[margin=23mm]{geometry}
\usepackage{fontspec}
\setmainfont{Latin Modern Roman}
\usepackage{amsmath,amssymb,mathtools}
\usepackage{xcolor,enumitem,booktabs,longtable,array,xurl,hyperref}
\definecolor{muted}{HTML}{53636C}
\hypersetup{colorlinks=true,urlcolor=blue,linkcolor=blue}
\setlength{\parindent}{0pt}
\setlength{\parskip}{5pt}
\newcommand{\R}{\mathbb R}
\newcommand{\E}{\mathbb E}
\newcommand{\N}{\mathbb N}
\newcommand{\ind}{\mathbf 1}
\newcommand{\Var}{\operatorname{Var}}
\newcommand{\Cov}{\operatorname{Cov}}
\newcommand{\supp}{\operatorname{supp}}
\DeclareMathOperator*{\argmin}{arg\,min}
\DeclareMathOperator*{\argmax}{arg\,max}
\newcommand{\entrymeta}[3]{{\sffamily\small\color{muted}\textbf{\texttt{#1}}\quad|\quad #2\par #3\par}}
\newcommand{\Problemref}[1]{\texttt{#1}}
\newcommand{\JELref}[2]{\texttt{#2} (JEL \texttt{#1})}
\begin{document}
\section*{Long-run reshoring effects}
\textbf{Problem ID:} \texttt{F13-409}\quad \textbf{JEL:} \texttt{F13}\par
% BEGIN ECONOMICS STATEMENT
\label{problem:OP-0343}
{\sffamily\small\color{muted}Original subject group: Trade, globalization and geoeconomics\par}
\label{definitions:problem:30007625}
\textbf{Audit classification:} Published research agenda. The 2026 primary draft explicitly asks about long-term capital/labor/productivity effects; the official agenda supports reconfigured trade. Capacity, persistence and rerouting definitions are editorial.
\subsection*{Definitions and precise research target}
Fix industries \(j\), countries and a tariff or sourcing-restriction policy path \(a\). Let \(K^d_{j,t+h}(a)\) be productive domestic capacity, defined by a prespecified engineering or production metric rather than nominal sales. Let \(V^d_{j,t+h}(a)\) be real domestic value added and \(C^{down}_{t+h}(a)\) downstream production costs. Define reshoring effects relative to an otherwise specified baseline path \(0\):
\[\Delta K_j(h)=E[K^d_{j,t+h}(do(a))-K^d_{j,t+h}(do(0))].\]
The task is to identify when positive capacity and value-added changes persist after the initial policy shock, alongside employment, productivity and downstream-cost effects. Specify the long horizon \(h\), policy permanence and capacity utilization. Distinguish relocation to the home country from import diversion through third countries, transshipment or increased imported inputs embedded in domestic assembly. Account for retaliation, foreign investment, exchange rates and endogenous entry. Identify or bound counterfactual capacity when policy selects industries already expected to expand or decline. Compare local adjustment estimates with national general-equilibrium outcomes. Evidence of reduced direct imports from a targeted partner does not establish durable domestic productive-capacity growth or a positive net welfare effect.

\textbf{Additional formal definitions and scope (editorial).} A policy path specifies its announcement, implementation, expected duration and any reversal date. Productive capacity is an engineering maximum at a stated utilization convention; realized output is capacity times endogenous utilization and productivity. Domestic value added subtracts imported and domestic intermediate purchases consistently and is deflated with a declared price index. Long-run persistence means a specified finite horizon or a limit \(\liminf_{h\to\infty}\Delta K_j(h)>0\), with stationary or balanced-growth normalization required before taking the limit. Define relocation by establishments or production stages, and identify ultimate country of origin to separate transshipment. A tariff-induced investment increase in protected firms may coexist with downstream contractions; national welfare requires all sectors and fiscal financing in the same equilibrium.

For the identification part, declare an admissible class \(\Theta\) of structural models, including preferences, technologies, shock laws, measurement/sampling rules and any equilibrium selector. If \(O\) denotes the complete observable history and \(T(\theta)\) the target defined above, the identified set is \(\mathcal I_T=\{T(\theta):\theta\in\Theta,\ \mathcal L_\theta(O)=\mathcal L_{\rm obs}(O)\}\); point identification means this set is a singleton. A \(do(a)\) comparison replaces stated policy/technology equations while fixing the declared exogenous law and re-solving the remaining equations. These definitions do not assert that an identification theorem holds without further restrictions.
\subsection*{Variants and assumption changes}
\begin{enumerate}
\item \textbf{Finite-horizon capacity response (editorial).} Fix policy permanence and a measurable capacity definition, then compare domestic investment with rerouted imports.
\item \textbf{Persistent general-equilibrium reshoring (source-motivated extension).} Include entry, downstream costs and foreign investment and use a long-horizon or balanced-growth limit; evaluate national welfare alongside capacity.
\end{enumerate}
\subsection*{References and source audit}
\begin{enumerate}
\item Official January 9, 2026 web version; locator: Interconnected World / Reconfigured trade. Broad agenda; exact model editorial. URL: \url{https://www.bankofengland.co.uk/research/bank-of-england-agenda-for-research}.
\item Brookings cover verifies Fajgelbaum and Khandelwal, March 2026 conference draft; NBER indexed metadata independently verifies WP 35064 and DOI. Locator: Section 7, printed p.49 (PDF p.50), explicit further-research list. Support: Dynamic and distributional tariff agenda. Full conference draft accessible. URL: \url{https://www.brookings.edu/wp-content/uploads/2026/03/1_Fajgelbaum-Khandelwal_unembargoed.pdf}.
\end{enumerate}
\begin{enumerate}\raggedright
\item\hypertarget{definitions:source:30007625:1}{} Bank of England. 2026. \emph{Bank of England Agenda for Research, 2025–2028}. Interconnected World / Reconfigured trade.
\url{https://www.bankofengland.co.uk/research/bank-of-england-agenda-for-research}.
\item\hypertarget{definitions:source:30007625:2}{} Pablo D. Fajgelbaum; Amit Khandelwal. 2026. \emph{Tariffs in 2025: Short-Run Impacts on the U.S. Economy}. Section 7, printed p.49 (PDF p.50), explicit further-research list.
\url{https://www.brookings.edu/wp-content/uploads/2026/03/1_Fajgelbaum-Khandelwal_unembargoed.pdf}.
DOI: \url{https://doi.org/10.3386/w35064}.
\end{enumerate}

\subsection*{Common conventions: Source mathematical conventions}

These conventions apply to each entry unless explicitly replaced.

\textbf{Models and identification.} A primitive model \(m\) specifies all probability laws, feasible choices, information, equilibrium and measurement rules used by its target. The admissible class \(\mathcal M\) is fixed independently of outcomes. If \(P_O(m)\) is its observable probability law and \(\tau(m)\) the target, its identified set at observable law \(P\) is
\[
 \mathcal I_\tau(P)=\{\tau(m):m\in\mathcal M,\ P_O(m)=P\}.
\]
Point identification means that this set is a singleton. An empty set signals incompatibility of data and assumptions. Coordinate bounds need not describe the joint set. Bounds are sharp only if every retained target is attainable or arbitrarily approachable by compatible models. Finite bounds require explicit restrictions on unobserved quantities; unrestricted confounding, hidden wealth, extrapolation or arbitrary equilibrium selection can yield uninformative sets.

\textbf{Causal interventions.} A potential outcome \(Y_i(p)\) is the outcome under a complete policy regime \(p\), including timing, financing, information and market spillovers. The target population and units are fixed before treatment. With interference, use the complete assignment vector or a justified exposure mapping; individual no-interference notation alone cannot identify an economy-wide intervention. A difference of mean potential outcomes does not require the cross-world joint distribution. Conditioning on post-treatment migration, survival, enrollment or employment changes the estimand and needs separate justification.

\textbf{Statistical statements.} An estimator, confidence set, forecast or test specifies a sampling law, model class and loss/coverage criterion. Uniform \((1-\alpha)\)-coverage means \(\inf_{m\in\mathcal M}\Pr_m\{\tau(m)\in C_n\}\geq1-\alpha\), or an explicitly asymptotic version. The whole parameter space trivially has coverage, so informativeness requires a stated width, risk or power objective. A forecasting improvement is relative to a named benchmark, information set and evaluation population.

\textbf{Equilibrium and optimization.} An equilibrium comprises feasible plans, optimality, beliefs where needed, budget constraints and all market/resource clearing. The primitives or a stated benchmark define each correspondence \(\mathcal E(p;m)\). Nonemptiness is assumed or proved, never inferred just from compact policy sets. A selected-equilibrium target uses a declared selection rule; a robust target quantifies over all permitted equilibria. Use suprema or infima unless attainment is established. An \(\varepsilon\)-optimal policy is feasible and within \(\varepsilon\) of the finite optimum value. Report an empty feasible set separately.

\textbf{Welfare and accounting.} Normative utility, interpersonal normalization, social weights, numeraire, discounting and the use of public receipts are primitives. Behavioral utility need not coincide with the welfare criterion. Household welfare includes ownership income and rents once; adding profits again to compensating variation generally double-counts them. A monetary equivalent is defined by an explicit transfer and price convention, with monotonicity and range conditions. Average effects, mechanism contributions, transfers and welfare losses are different targets. Infinite sums require convergence and dynamic feasible plans require terminal or no-Ponzi conditions.

\textbf{Source versions.} A working-paper date, revision date and journal publication year are distinguished. A DOI for a journal version is not automatically the DOI of an earlier manuscript. Locators refer to the stated version and printed pages where specified. A metadata-only check does not verify a claim about a full-text conclusion. Every entry's source subsection narrows the provenance claim accordingly.
% END ECONOMICS STATEMENT
\end{document}
